\documentclass[12pt]{article}

\usepackage[spanish]{babel}
\usepackage[utf8]{inputenc}
\usepackage[T1]{fontenc}
\usepackage{amsmath,amssymb,amsfonts}
\usepackage{physics}
\usepackage{geometry}
\usepackage{booktabs}
\usepackage{bm}

\geometry{margin=2.5cm}

\title{Sistemas Coordenados Ortogonales y su Relación con la Formulación Tensorial}
\author{Matemáticas Avanzadas de la Física}
\date{}

\begin{document}
\fontsize{14}{14}\selectfont
\maketitle

\section{Introducción.}

En coordenadas cartesianas, las expresiones para gradiente, divergencia y rotacional poseen una forma particularmente simple debido a que:
\begin{itemize}
    \item la métrica es constante,
    \item los vectores base no dependen de la posición,
    \item los factores de escala son unitarios.
\end{itemize}
Al cambiar a coordenadas cilíndricas, esféricas u otros sistemas curvilíneos, estas propiedades dejan de cumplirse y la geometría del espacio aparece explícitamente en las ecuaciones.
\par<
El objetivo de este documento es mostrar cómo los conceptos de:
\begin{itemize}
    \item reglas de transformación,
    \item factores de escala,
    \item operadores diferenciales,
\end{itemize}
pueden entenderse naturalmente mediante:
\begin{itemize}
    \item notación de índices,
    \item tensores,
    \item símbolo de Levi-Civita,
    \item símbolos de Christoffel.
\end{itemize}

\section{Transformaciones de coordenadas.}

Consideremos una transformación entre coordenadas cartesianas y un sistema general:
\begin{align*}
    x^{i} = x^{i} \left( q^{1}, q^{2}, q^{3} \right)
\end{align*}
Por ejemplo, para coordenadas esféricas:
\begin{align*}
x &= r \sin \theta \cos \phi \\
y &= r \sin \theta \sin \phi \\
z &= r \cos \theta
\end{align*}
El Jacobiano de la transformación es:
\begin{align*}
    J^{i}_{j} = \pdv{x^{i}}{q^{j}}
\end{align*}
La transformación de un vector contravariante está dada por:
\begin{align*}
    V^{\prime \, {i}} = \pdv{q^{\prime \, i}}{q^{j}} \, V^{j}
\end{align*}
mientras que un covector transforma según:

\[
V'_i=
\frac{\partial q^j}{\partial q'^i}V_j
\]

Estas leyes motivan la definición de tensor.

\section{Notación de índices y convención de Einstein.}

Se adopta la convención de suma de Einstein:

\[
A_iB^i=
\sum_{i=1}^{3}A_iB^i
\]

Por ejemplo, el producto punto se escribe simplemente como:

\[
\mathbf{A}\cdot\mathbf{B}=A_iB^i
\]

Un tensor de segundo orden transforma como:

\[
T'^i_{\ j}
=
\frac{\partial q'^i}{\partial q^k}
\frac{\partial q^l}{\partial q'^j}
T^k_{\ l}
\]

\section{Factores de escala y tensor métrico}

En coordenadas ortogonales el elemento diferencial de longitud puede escribirse como:

\[
d\mathbf{r}
=
h_1dq^1\hat{e}_1+
h_2dq^2\hat{e}_2+
h_3dq^3\hat{e}_3
\]

donde

\[
h_i=
\left|
\frac{\partial \mathbf{r}}
{\partial q^i}
\right|
\]

son los factores de escala.

El elemento de línea resulta:

\[
ds^2=
h_1^2(dq^1)^2+
h_2^2(dq^2)^2+
h_3^2(dq^3)^2
\]

que puede escribirse en forma tensorial como:

\[
ds^2=
g_{ij}dq^idq^j
\]

En coordenadas ortogonales:

\[
g_{ij}=0
\qquad
(i\neq j)
\]

y

\[
g_{ii}=h_i^2
\]

por lo que:

\[
h_i=\sqrt{g_{ii}}
\]

Por tanto, los factores de escala corresponden a las raíces cuadradas de las componentes diagonales del tensor métrico.

\section{Operadores diferenciales mediante índices}

\subsection{Gradiente}

El gradiente de un escalar puede escribirse como:

\[
(\nabla \phi)_i=
\partial_i\phi
\]

\subsection{Divergencia}

La divergencia de un vector es:

\[
\nabla\cdot\mathbf{A}
=
\partial_iA^i
\]

\subsection{Rotacional}

El rotacional se expresa utilizando el símbolo de Levi-Civita:

\[
(\nabla\times\mathbf{A})^i
=
\varepsilon^{ijk}
\partial_jA_k
\]

\section{Símbolo de Levi-Civita}

El símbolo de Levi-Civita se define mediante:

\[
\varepsilon_{123}=1
\]

y cambia de signo bajo permutaciones impares:

\[
\varepsilon_{213}=-1
\]

mientras que vale cero si existen índices repetidos.

El producto cruz puede escribirse como:

\[
(\mathbf{A}\times\mathbf{B})^i
=
\varepsilon^{ijk}
A_jB_k
\]

lo que permite unificar el tratamiento del producto vectorial y del rotacional.

\section{Determinante métrico}

En sistemas curvilíneos aparece el determinante de la métrica:

\[
g=\det(g_{ij})
\]

y el tensor completamente antisimétrico se define como:

\[
E^{ijk}
=
\frac{1}{\sqrt{g}}
\varepsilon^{ijk}
\]

En coordenadas esféricas:

\[
g=r^4\sin^2\theta
\]

por lo tanto:

\[
\sqrt{g}=r^2\sin\theta
\]

Esto explica la aparición de factores geométricos en los operadores diferenciales expresados en coordenadas esféricas.

\section{Dependencia espacial de las bases}

En coordenadas cartesianas:

\[
\frac{\partial \hat{e}_i}{\partial x^j}=0
\]

Sin embargo, en coordenadas esféricas:

\[
\frac{\partial \hat{r}}{\partial\theta}
=
\hat{\theta}
\]

y

\[
\frac{\partial \hat{r}}{\partial\phi}
=
\sin\theta\hat{\phi}
\]

Los vectores base dependen de la posición y sus derivadas ya no son nulas.

\section{Símbolos de Christoffel}

Debido a la variación espacial de las bases, la derivada parcial de un vector no transforma tensorialmente.

Para corregir esto se introduce la derivada covariante:

\[
\nabla_jV^i
=
\partial_jV^i
+
\Gamma^i_{\ jk}V^k
\]

donde:

\[
\Gamma^i_{\ jk}
=
\frac{1}{2}
g^{im}
\left(
\partial_jg_{mk}
+
\partial_kg_{mj}
-
\partial_mg_{jk}
\right)
\]

Los símbolos de Christoffel representan el cambio geométrico de la base local.

\section{Divergencia mediante derivada covariante}

La divergencia puede escribirse como:

\[
\nabla_iV^i
=
\partial_iV^i
+
\Gamma^i_{\ ik}V^k
\]

o equivalentemente:

\[
\nabla_iV^i
=
\frac{1}{\sqrt{g}}
\partial_i
\left(
\sqrt{g}V^i
\right)
\]

En coordenadas ortogonales:

\[
\sqrt{g}=h_1h_2h_3
\]

\section{Ejemplo: coordenadas cilíndricas}

Los factores de escala son:

\[
h_r=1,
\qquad
h_{\phi}=r,
\qquad
h_z=1
\]

Por tanto:

\[
\sqrt{g}=r
\]

La divergencia toma la forma:

\[
\nabla\cdot\mathbf{A}
=
\frac{1}{r}
\frac{\partial}{\partial r}
(rA_r)
+
\frac{1}{r}
\frac{\partial A_{\phi}}{\partial \phi}
+
\frac{\partial A_z}{\partial z}
\]

sin necesidad de memorizar la expresión particular.

\section{Secuencia didáctica sugerida}

\begin{enumerate}
    \item Transformaciones de coordenadas.
    \item Jacobiano y regla de la cadena.
    \item Convención de Einstein.
    \item Vectores covariantes y contravariantes.
    \item Factores de escala.
    \item Tensor métrico.
    \item Gradiente, divergencia y rotacional.
    \item Símbolo de Levi-Civita.
    \item Dependencia espacial de las bases.
    \item Derivada covariante.
    \item Símbolos de Christoffel.
    \item Operadores diferenciales en forma tensorial.
\end{enumerate}

\section{Conclusión}

La formulación tensorial permite unificar el tratamiento de todos los sistemas coordenados.

\begin{center}
\begin{tabular}{ll}
\toprule
Concepto clásico & Interpretación geométrica \\
\midrule
Factores de escala & Componentes de la métrica \\
Jacobiano & Transformación tensorial \\
Producto cruz & Levi-Civita \\
Cambio de versores & Christoffel \\
Gradiente, divergencia y rotacional & Derivadas covariantes \\
Sistemas ortogonales & Caso particular de geometría riemanniana \\
\bottomrule
\end{tabular}
\end{center}

De esta forma, las expresiones particulares de los operadores diferenciales dejan de ser fórmulas aisladas para memorizar y pasan a entenderse como manifestaciones de una estructura geométrica única.


\section{Velocidad y Aceleración mediante Símbolos de Christoffel.}

Sea una trayectoria parametrizada mediante coordenadas generalizadas
\begin{align*}
    q^{i} = q^{i} (t)
\end{align*}
La velocidad contravariante se define como
\begin{align*}
    v^{i} = \dot{q}^{i}
\end{align*}
mientras que la aceleración geométrica se obtiene mediante la derivada covariante de la velocidad:
\begin{align*}
    a^{i} = \pdv{v^{i}}{t}
\end{align*}
es decir,
\begin{align*}
    \boxed{
        a^{i} = \ddot{q}^{i} + \Gamma_{jk}^{i} \, \dot{q}^{j} \, \dot{q}^{k}
    }
\end{align*}
Los términos que contienen símbolos de Christoffel representan la contribución asociada a la variación espacial de la base local.

%%%%%%%%%%%%%%%%%%%%%%%%%%%%%%%%%%%%%%%%%%%%%%%%%%%%%%%%%%%
\section{Coordenadas esféricas.}

Las coordenadas son

\[
q^1=r,
\qquad
q^2=\theta,
\qquad
q^3=\phi
\]

con transformación

\begin{align}
x&=r\sin\theta\cos\phi\\
y&=r\sin\theta\sin\phi\\
z&=r\cos\theta
\end{align}

\subsection{Tensor métrico}

El elemento diferencial de longitud es

\[
ds^2=
dr^2+
r^2d\theta^2+
r^2\sin^2\theta\,d\phi^2
\]

por lo tanto,

\[
g_{ij}
=
\begin{pmatrix}
1&0&0\\
0&r^2&0\\
0&0&r^2\sin^2\theta
\end{pmatrix}
\]

y

\[
g^{ij}
=
\begin{pmatrix}
1&0&0\\
0&\frac1{r^2}&0\\
0&0&
\frac1{r^2\sin^2\theta}
\end{pmatrix}
\]

\subsection{Símbolos de Christoffel no nulos}

Utilizando

\[
\Gamma^i_{jk}
=
\frac12
g^{im}
\left(
\partial_jg_{mk}
+
\partial_kg_{mj}
-
\partial_mg_{jk}
\right)
\]

se obtiene:

\begin{align}
\Gamma^r_{\theta\theta}
&=-r
\\
\Gamma^r_{\phi\phi}
&=
-r\sin^2\theta
\\
\Gamma^\theta_{r\theta}
=
\Gamma^\theta_{\theta r}
&=
\frac1r
\\
\Gamma^\theta_{\phi\phi}
&=
-\sin\theta\cos\theta
\\
\Gamma^\phi_{r\phi}
=
\Gamma^\phi_{\phi r}
&=
\frac1r
\\
\Gamma^\phi_{\theta\phi}
=
\Gamma^\phi_{\phi\theta}
&=
\cot\theta
\end{align}

\subsection{Velocidad}

Las componentes contravariantes de la velocidad son

\[
v^r=\dot r
\]

\[
v^\theta=\dot\theta
\]

\[
v^\phi=\dot\phi
\]

Las componentes físicas se obtienen multiplicando por los factores de escala:

\[
h_r=1,
\qquad
h_\theta=r,
\qquad
h_\phi=r\sin\theta
\]

Por tanto,

\[
\boxed{
\mathbf v
=
\dot r\,\hat r
+
r\dot\theta\,\hat\theta
+
r\sin\theta\,\dot\phi\,\hat\phi
}
\]

\subsection{Aceleración}

\subsubsection*{Componente radial}

\[
a^r
=
\ddot r
+
\Gamma^r_{\theta\theta}\dot\theta^2
+
\Gamma^r_{\phi\phi}\dot\phi^2
\]

\[
a^r
=
\boxed{
\ddot r
-r\dot\theta^2
-r\sin^2\theta\dot\phi^2
}
\]

\subsubsection*{Componente polar}

\[
a^\theta
=
\ddot\theta
+
2\Gamma^\theta_{r\theta}
\dot r\dot\theta
+
\Gamma^\theta_{\phi\phi}\dot\phi^2
\]

\[
a^\theta
=
\boxed{
\ddot\theta
+
\frac{2}{r}\dot r\dot\theta
-
\sin\theta\cos\theta\dot\phi^2
}
\]

\subsubsection*{Componente azimutal}

\[
a^\phi
=
\ddot\phi
+
2\Gamma^\phi_{r\phi}\dot r\dot\phi
+
2\Gamma^\phi_{\theta\phi}
\dot\theta\dot\phi
\]

\[
a^\phi
=
\boxed{
\ddot\phi
+
\frac{2}{r}\dot r\dot\phi
+
2\cot\theta\dot\theta\dot\phi
}
\]

Las componentes físicas son

\[
\boxed{
\mathbf a=
a_r\hat r
+
a_\theta\hat\theta
+
a_\phi\hat\phi
}
\]

con

\begin{align}
a_r&=
\ddot r
-r\dot\theta^2
-r\sin^2\theta\dot\phi^2
\\
a_\theta&=
r\ddot\theta
+2\dot r\dot\theta
-r\sin\theta\cos\theta\dot\phi^2
\\
a_\phi&=
r\sin\theta\ddot\phi
+2\sin\theta\dot r\dot\phi
+2r\cos\theta\dot\theta\dot\phi
\end{align}

%%%%%%%%%%%%%%%%%%%%%%%%%%%%%%%%%%%%%%%%%%%%%%%%%%%%%%%%%%%
\section{Coordenadas cilíndricas parabólicas}

Las coordenadas son

\[
q^1=u,
\qquad
q^2=v,
\qquad
q^3=z
\]

definidas mediante

\begin{align}
x&=
\frac12(u^2-v^2)
\\
y&=
uv
\\
z&=z
\end{align}

\subsection{Factores de escala}

Los vectores tangentes son

\[
\frac{\partial\mathbf r}{\partial u}
=
(u,v,0)
\]

\[
\frac{\partial\mathbf r}{\partial v}
=
(-v,u,0)
\]

por lo que

\[
h_u=h_v=
\sqrt{u^2+v^2}
\]

y

\[
h_z=1
\]

El elemento diferencial de longitud es

\[
ds^2=
(u^2+v^2)
(du^2+dv^2)
+dz^2
\]

Por tanto,

\[
g_{ij}
=
\begin{pmatrix}
u^2+v^2&0&0\\
0&u^2+v^2&0\\
0&0&1
\end{pmatrix}
\]

\subsection{Métrica inversa}

\[
g^{ij}
=
\begin{pmatrix}
\dfrac1{u^2+v^2}&0&0\\
0&
\dfrac1{u^2+v^2}&0\\
0&0&1
\end{pmatrix}
\]

\subsection{Símbolos de Christoffel no nulos}

Los símbolos distintos de cero son

\begin{align}
\Gamma^u_{uu}
&=
\frac{u}{u^2+v^2}
\\
\Gamma^u_{uv}
=
\Gamma^u_{vu}
&=
\frac{v}{u^2+v^2}
\\
\Gamma^u_{vv}
&=
-\frac{u}{u^2+v^2}
\\
\Gamma^v_{uu}
&=
-\frac{v}{u^2+v^2}
\\
\Gamma^v_{uv}
=
\Gamma^v_{vu}
&=
\frac{u}{u^2+v^2}
\\
\Gamma^v_{vv}
&=
\frac{v}{u^2+v^2}
\end{align}

\subsection{Velocidad}

Las componentes físicas de la velocidad son

\[
\boxed{
\mathbf v
=
\sqrt{u^2+v^2}\,
\dot u\,\hat e_u
+
\sqrt{u^2+v^2}\,
\dot v\,\hat e_v
+
\dot z\,\hat z
}
\]

\subsection{Aceleración}

\subsubsection*{Componente \(u\)}

\[
a^u=
\ddot u
+
\Gamma^u_{uu}\dot u^2
+
2\Gamma^u_{uv}\dot u\dot v
+
\Gamma^u_{vv}\dot v^2
\]

\[
\boxed{
a^u=
\ddot u
+
\frac{
u(\dot u^2-\dot v^2)
+
2v\dot u\dot v
}
{u^2+v^2}
}
\]

\subsubsection*{Componente \(v\)}

\[
a^v=
\ddot v
+
\Gamma^v_{uu}\dot u^2
+
2\Gamma^v_{uv}\dot u\dot v
+
\Gamma^v_{vv}\dot v^2
\]

\[
\boxed{
a^v=
\ddot v
+
\frac{
-v\dot u^2
+
2u\dot u\dot v
+
v\dot v^2
}
{u^2+v^2}
}
\]

\subsubsection*{Componente \(z\)}

\[
a^z=\ddot z
\]

Las componentes físicas se obtienen multiplicando por los factores de escala correspondientes.

%%%%%%%%%%%%%%%%%%%%%%%%%%%%%%%%%%%%%%%%%%%%%%%%%%%%%%%%%%%

\section{Conclusiones}

La expresión

\[
a^i=
\ddot q^i+
\Gamma^i_{jk}
\dot q^j
\dot q^k
\]

permite obtener de manera sistemática las aceleraciones en cualquier sistema coordenado sin necesidad de derivar explícitamente los versores unitarios.

Los símbolos de Christoffel representan la contribución asociada al cambio espacial de la base local y constituyen la generalización geométrica de las derivadas de los vectores unitarios utilizadas tradicionalmente en mecánica clásica.


\end{document}